% ============================================================================
% Fichier  : partie2/P02_C06_analyse_reseau.tex
% Titre    : L'Analyse du Réseau
% Auteur   : MINKA MI NGUIDJOÏ Thierry Emmanuel
% Version  : 1.0 -- novembre 2026
% Corrigés : corriges/Corr_C06.tex
% ============================================================================
% Sources : Bejtlich (2013) pour la supervision de la sécurité réseau ;
% Sanders (2017) pour l'analyse de paquets ; Wireshark, tshark, Zeek
% (présentés sans citation) ; RFC 5737 (adresses de documentation).
% Faits nouveaux (bible) : adresses internes SRV-AD01 10.20.1.10, SRV-FIC01
% 10.20.1.20, SRV-CPT01 10.20.1.30 ; lien Internet de 20 Mbit/s ; flux
% SRV-FIC01 -> 203.0.113.47:443 du 14/03 de 1 h 12 à 3 h 47, 3,2 Go montant,
% 11 Mo descendant, SNI cdn-update.example.net, certificat auto-signé du
% 01/03/2025 ; balise de SRV-AD01 toutes les 60 s (+/- 6 s) ; virements
% émis depuis 203.0.113.62, hors du réseau de NYANGA ; journal de l'outil de
% copie sur SRV-FIC01 : 18 412 fichiers du dossier Clients.
% ============================================================================

\chapter{L'Analyse du Réseau}
\label{chap:reseau}
\chaptermark{L'analyse du réseau}

\epigraph{%
    « Le chiffrement cache le contenu.\\
    Il ne cache ni la taille, ni l'heure, ni la destination. »%
}{Principe du cours}

\niveauChapitre{Intermédiaire}{%
    Maîtrise des Chapitres~\ref{chap:acquisition} et~\ref{chap:memoire}
}

\begin{boxobjectifs}
\begin{itemize}
    \item Distinguer les sources réseau~: capture complète, flux, journaux
          d'équipements, alertes de sonde.
    \item Lire une capture avec un analyseur de paquets et en extraire les
          conversations utiles.
    \item Caractériser une exfiltration par le volume, la durée, le débit
          et la destination.
    \item Reconnaître une balise de commande par sa régularité.
    \item Tirer d'un flux chiffré tout ce qu'il dit sans être déchiffré.
    \item Dire ce que l'absence d'un flux dans une capture établit, et ce
          qu'elle n'établit pas.
\end{itemize}
\end{boxobjectifs}

\bigskip

% ============================================================================
\section*{Avant de commencer : 48 heures de trafic}
% ============================================================================

Yaoundé, 24 mars 2025. La capture de la sonde de NYANGA CRÉDIT (scellé
NC-04) couvre le lien Internet du siège du 13~mars à 8~h au 14~mars à 8~h~:
37~Go de paquets. M.~FOTSO, analyste de SANKOFA EXPERTISE, a trois questions
de Mme~MBARGA~: qu'est-ce qui est sorti, vers où, et à quel moment~; le
programme du contrôleur de domaine parlait-il à la même adresse~; les
virements de la nuit sont-ils passés par ce lien~?

\begin{boxalert}
\textbf{Violation typique.} Conclure qu'un échange n'a pas eu lieu parce
qu'on ne le trouve pas dans la capture.

\smallskip
\textbf{Règle~:} une capture ne voit que ce qui passe par le point où elle
est placée, pendant la période qu'elle couvre. Avant d'écrire « aucun flux
vers telle adresse », on dit où était la sonde, ce qu'elle voyait, et si le
flux pouvait emprunter un autre chemin.
\end{boxalert}

\separateur

% ============================================================================
\section{Quatre Sources Réseau}
\label{sec:c06-sources}
% ============================================================================

\begin{tableaucours}{L{2.8cm}L{4.0cm}L{2.6cm}L{2.8cm}}{%
    \tableauHeader{Source} &
    \tableauHeader{Contenu} &
    \tableauHeader{Durée typique} &
    \tableauHeader{Limite}
}
Capture complète (PCAP) & Tous les paquets, en-têtes et contenu & Heures à
jours (volume) & Coût de stockage~; contenu souvent chiffré \\
Flux (NetFlow, IPFIX) & Qui a parlé à qui, quand, combien d'octets & Mois &
Pas de contenu \\
Journaux d'équipements & Pare-feu, proxy, DNS~: décisions et requêtes &
Semaines à mois & Ce que l'équipement a jugé utile de noter \\
Alertes de sonde & Événements jugés suspects par des règles & Mois & Ne
voit que ce que les règles décrivent \\
\end{tableaucours}
\vspace{-0.8em}
\begin{center}{\small\itshape Tableau~6.1. Les sources de l'analyse réseau}\end{center}

R.~Bejtlich, dans son ouvrage sur la supervision de la sécurité
réseau~\cite{bejtlich2013nsm}, montre que ces sources se complètent~: les
flux et les journaux couvrent une longue période sans contenu, la capture
complète couvre une période courte avec tout le détail. Une enquête commence
souvent par les flux, pour repérer la fenêtre utile, puis descend dans la
capture.

% ============================================================================
\section{Lire une Capture}
\label{sec:c06-lire}
% ============================================================================

L'analyseur de paquets Wireshark, et sa version en ligne de commande tshark,
décodent les protocoles et regroupent les paquets en
\termecle{conversations} (une paire d'adresses et de ports). Les outils de
supervision comme Zeek transforment une capture en journaux par protocole~:
connexions, requêtes DNS, échanges TLS, transferts HTTP. On travaille dans
cet ordre.

\begin{enumerate}
    \item Les statistiques de conversations~: quelles paires ont échangé le
          plus d'octets, dans quel sens.
    \item Les requêtes DNS~: quels noms ont été résolus, vers quelles
          adresses.
    \item Les échanges TLS~: nom de serveur annoncé (SNI), certificat
          présenté, versions.
    \item Le détail des paquets, seulement pour les conversations retenues.
\end{enumerate}

\begin{boxinfo}
\textbf{Toujours sur une copie, toujours avec un filtre écrit.} On ouvre la
copie de travail de la capture, jamais le scellé, et l'on consigne chaque
filtre appliqué (par exemple \texttt{ip.addr == 203.0.113.47}). Le filtre
fait partie de la méthode~: un autre analyste doit pouvoir le rejouer.
\end{boxinfo}

% ============================================================================
\section{Caractériser une Exfiltration}
\label{sec:c06-exfiltration}
% ============================================================================

Une exfiltration se reconnaît moins à son contenu, souvent chiffré, qu'à sa
forme~: un volume sortant important, un rapport entre octets envoyés et
reçus très déséquilibré, une destination inhabituelle, un horaire anormal,
un débit soutenu.

\begin{boxscenario}{Opération NYANGA. La conversation de 1 h 12}
\textbf{Constatations} (copie de NC-04).
\begin{itemize}
    \item Le 14~mars, de 1~h~12 à 3~h~47, une conversation TLS relie
          SRV-FIC01 (10.20.1.20) à 203.0.113.47, port~443.
    \item Octets envoyés par SRV-FIC01~: 3,2~Go~; octets reçus~: 11~Mo.
    \item Le client annonce le nom \texttt{cdn-update.example.net}~; le
          serveur présente un certificat auto-signé émis le 1\textsuperscript{er}~mars
          2025.
    \item Aucune autre conversation de SRV-FIC01 vers Internet sur les
          24~heures de la capture.
\end{itemize}

\textbf{Calcul.} Durée~: 2~h~35, soit 9\,300~s. Débit moyen~: $3{,}2 \times
10^{9} / 9\,300 \approx 344\,000$~octets par seconde, soit environ
2,75~Mbit/s, sur un lien de 20~Mbit/s. Rapport envoyé sur reçu~: environ
290.

\textbf{Analyse} (confiance élevée). Un serveur de fichiers n'envoie pas
3,2~Go en pleine nuit vers un serveur au certificat auto-signé récent qui
se fait passer pour un service de mise à jour. La forme du flux est celle
d'une exfiltration. Le débit, limité à un septième du lien, est compatible
avec un outil réglé pour ne pas saturer la ligne et passer inaperçu.
\end{boxscenario}

\subsection{Ce que le chiffrement laisse voir}

\begin{tableaucours}{L{4.2cm}L{8.0cm}}{%
    \tableauHeader{Élément visible} &
    \tableauHeader{Ce qu'il apprend}
}
Adresses, ports, heures & Qui, vers où, quand \\
Volumes dans chaque sens & Envoi ou réception, ordre de grandeur des
données \\
Nom annoncé (SNI) & Le service que le client prétend joindre \\
Certificat du serveur & Émetteur, date de création, auto-signature \\
Rythme des paquets & Transfert continu, balise régulière, session
interactive \\
\end{tableaucours}
\vspace{-0.8em}
\begin{center}{\small\itshape Tableau~6.2. Ce que laisse voir un flux chiffré}\end{center}

Le contenu, lui, reste inaccessible sans les clés. Pour savoir quels
fichiers sont partis, il faut une autre source. À NYANGA, elle se trouve sur
SRV-FIC01 lui-même~: l'outil de copie utilisé par l'attaquant a laissé un
journal dans le dossier de l'utilisateur, retrouvé dans l'acquisition du
18~mars (NC-07). Il liste 18\,412 fichiers du dossier \texttt{Clients},
pour un total de 3,19~Go, avec des heures de début et de fin qui encadrent
la conversation réseau.

% ============================================================================
\section{Reconnaître une Balise}
\label{sec:c06-balise}
% ============================================================================

\begin{boxdef}
\textbf{\termecle{Balise}} (\textit{beacon})~: connexion qu'un programme
malveillant établit à intervalles réguliers vers son serveur de commande,
pour signaler sa présence et recevoir des ordres. Sa régularité la trahit,
même quand le contenu est chiffré et le volume faible.
\end{boxdef}

\begin{boxscenario}{Opération NYANGA. Toutes les soixante secondes}
\textbf{Constatations.} Sur les 24~heures de la capture, SRV-AD01
(10.20.1.10) ouvre 1\,431 connexions courtes vers 203.0.113.47:443.
L'intervalle entre deux connexions est compris entre 54 et 66~secondes~;
chaque échange fait entre 1,2 et 2,8~Ko, sauf onze échanges de 40 à
900~Ko, tous situés entre 0~h~50 et 3~h~50 le 14~mars.

\textbf{Calcul.} 24~heures font 1\,440 minutes~: 1\,431 connexions
correspondent à une connexion par minute, à neuf près.

\textbf{Analyse} (confiance élevée). C'est la balise du PID~4412 vu en
mémoire (Chapitre~\ref{chap:memoire}), réglée sur 60~secondes avec une
variation de 10~\%. Les onze échanges plus volumineux, pendant la nuit,
correspondent à des ordres plus lourds~: ils précèdent et accompagnent
l'exfiltration depuis SRV-FIC01.
\end{boxscenario}

% ============================================================================
\section{Ce que la Capture ne Voit pas}
\label{sec:c06-absence}
% ============================================================================

La troisième question de Mme~MBARGA reçoit une réponse inattendue. La
capture ne contient aucun échange avec l'interface de paiement pendant la
nuit. Les journaux de l'interface (NC-05) indiquent que les vingt-trois
virements ont été émis depuis l'adresse 203.0.113.62, qui n'appartient pas à
NYANGA CRÉDIT. L'interface est hébergée par le prestataire, sur Internet~:
l'attaquant a utilisé la clé volée depuis sa propre machine, sans passer par
le réseau du siège.

\begin{boxlecture}{Les trois questions au réseau}
\begin{tableaucours}{L{7.0cm}cL{3.4cm}}{%
    \tableauHeader{Proposition} &
    \tableauHeader{Cote} &
    \tableauHeader{Sources}
}
3,2~Go sont sortis de SRV-FIC01 vers 203.0.113.47 la nuit du 14~mars &
\OUI{} & NC-04 \\
Les fichiers sortis sont ceux du dossier \texttt{Clients} & \OUI{} &
Journal de l'outil (NC-07) \\
Le programme de SRV-AD01 communiquait avec la même adresse & \OUI{} & NC-04,
NC-01 \\
Les virements sont passés par le lien du siège & \NON{} & NC-04, NC-05 \\
Aucune autre donnée n'est sortie par un autre chemin & \NEK{} & La sonde ne
voit que le lien Internet du siège \\
\end{tableaucours}

\medskip
\textbf{Lecture.} Trois \OUI{}, un \NON{}, un \NEK{}. Le \NON{} sur les
virements est solide parce que les deux sources se recoupent~: absence dans
la capture, et adresse d'émission externe dans les journaux de l'interface.
Le \NEK{} rappelle que les agences ont leur propre accès Internet, hors de la
vue de la sonde.
\end{boxlecture}

% ============================================================================
\section{Ce que l'Analyse Réseau demande à chaque Métier}
\label{sec:c06-metiers}
% ============================================================================

\begin{boxmetier}{Analyste en réponse à incident}
\textbf{Analyser une capture}
\begin{enumerate}
    \item Situer la sonde~: lien observé, période, pertes éventuelles.
    \item Partir des statistiques de conversations, pas des paquets.
    \item Pour chaque conversation suspecte~: volumes, durée, débit,
          rythme, nom annoncé, certificat.
    \item Rattacher chaque flux à un processus (mémoire) et à des fichiers
          (disque).
    \item Écrire les filtres utilisés.
\end{enumerate}
\end{boxmetier}

\begin{boxmetier}{RSSI}
Placer la sonde là où passent les flux qui comptent, et savoir ce qu'elle
ne voit pas. À NYANGA, les vingt-deux agences sortent sur Internet par leurs
propres accès~: elles sont hors du champ de la sonde.
\end{boxmetier}

\begin{boxmetier}{Application OPJ/Expert Judiciaire}
L'adresse 203.0.113.47 désigne une machine à un instant donné. Pour aller
vers une personne, l'enquête doit obtenir, par réquisition ou par entraide
internationale, les données du fournisseur qui l'exploitait~: le
Chapitre~\ref{chap:normes-droit} en donne les fondements.
\end{boxmetier}

\separateur

\begin{boxresume}
\textbf{Ce qu'il faut retenir de ce chapitre~:}
\begin{itemize}
    \item Capture complète, flux, journaux et alertes se complètent~; on
          part du large pour descendre vers le détail.
    \item Une exfiltration se caractérise par son volume, son déséquilibre,
          son débit, son horaire et sa destination~; une balise par sa
          régularité.
    \item Un flux chiffré laisse voir les adresses, les volumes, le nom
          annoncé, le certificat et le rythme~; son contenu se cherche
          ailleurs.
    \item L'absence d'un flux dans une capture ne vaut \NON{} que si l'on
          sait que le flux devait passer par la sonde.
    \item À NYANGA, 3,2~Go du dossier \texttt{Clients} sont partis de
          SRV-FIC01~; les virements, eux, ont été émis de l'extérieur.
\end{itemize}
\end{boxresume}

\begin{boxinfo}
\textbf{Points de contrôle.} Vous avez compris ce chapitre si vous savez~:
\begin{itemize}
    \item calculer le débit moyen d'une conversation~;
    \item reconnaître une balise dans une liste de connexions~;
    \item dire pourquoi l'absence d'un flux ne prouve pas toujours son
          inexistence.
\end{itemize}
\end{boxinfo}

% ============================================================================
\section*{Exercices}
% ============================================================================

\begin{boxexo}[title={\ding{45}\quad Exercice~6.1. Quelle source~? \niveaubase}]
Pour chaque question, dites quelle source réseau y répond le mieux~:
(a) SRV-FIC01 a-t-il parlé à 203.0.113.47 il y a trois mois~? (b) quel nom
de domaine le contrôleur de domaine a-t-il résolu avant sa première balise~?
(c) que contenait la requête HTTP non chiffrée d'un poste vers un site
inconnu~? (d) la sonde a-t-elle signalé quelque chose cette nuit-là~?
\end{boxexo}

\begin{boxexo}[title={\ding{45}\quad Exercice~6.2. Le calcul du débit \niveaumoyen}]
Une seconde conversation, découverte dans les flux d'une agence, montre
900~Mo envoyés en 1~h~15.
\begin{enumerate}[(a)]
    \item Calculez le débit moyen en octets par seconde et en Mbit/s.
    \item L'agence dispose d'un lien de 4~Mbit/s. Qu'en déduisez-vous sur le
          réglage de l'outil~?
    \item Cette conversation change-t-elle une cote de l'encadré « Les trois
          questions au réseau »~?
\end{enumerate}
\end{boxexo}

\begin{boxexo}[title={\ding{45}\quad Exercice~6.3. Une balise ou pas \niveaumoyen}]
Voici les heures de connexion d'un poste vers une même adresse~: 10:00:03,
10:05:01, 10:10:04, 10:15:02, 10:20:03, puis 11:47:30, 12:13:05.
\begin{enumerate}[(a)]
    \item Lesquelles évoquent une balise, et de quelle période~?
    \item Proposez une explication légitime de ce rythme.
    \item Quelles vérifications feriez-vous avant de conclure~?
\end{enumerate}
\end{boxexo}

\begin{boxexo}[title={\ding{45}\quad Exercice~6.4. Le NON des virements \niveauexpert}]
\begin{enumerate}[(a)]
    \item Expliquez pourquoi la seule absence des virements dans la capture
          ne suffirait pas à coter \NON{} la proposition « les virements
          sont passés par le lien du siège ».
    \item Quelle pièce indépendante permet ce \NON{}~? Quelle est sa propre
          faiblesse (Chapitre~\ref{chap:preuve})~?
    \item Proposez une démarche pour renforcer la cote, en vous plaçant du
          côté de l'enquête judiciaire.
\end{enumerate}
\end{boxexo}

\bigskip

\begin{boxinfo}
\textbf{Pour aller plus loin.}
\begin{itemize}
    \item R.~Bejtlich, \textit{The Practice of Network Security
          Monitoring}~\cite{bejtlich2013nsm}.
    \item C.~Sanders, \textit{Practical Packet Analysis}~\cite{sanders2017packet}.
\end{itemize}
\end{boxinfo}

\begin{boxinfo}
\textbf{Transition.} Le disque, la mémoire et le réseau ont chacun parlé,
chacun avec ses heures. Il reste à les mettre dans un même ordre, sur une
même échelle de temps~: Chapitre~\ref{chap:chronologie}.
\end{boxinfo}
